\documentclass[11pt]{article}
\usepackage[a4paper,margin=27mm]{geometry}
\usepackage{amsmath,amssymb,mathtools}
\usepackage{fontspec,unicode-math}
\setmainfont{Latin Modern Roman}
\setmathfont{Latin Modern Math}
\usepackage{xurl}
\usepackage[hidelinks]{hyperref}
\newcommand{\sref}[2]{\hyperlink{source:#1:#2}{[S#2]}}
\newcommand{\cataloguescope}[1]{\paragraph{Mathematical question.}#1}
\hypersetup{pdftitle={Problem statement}}
\begin{document}
\section{Wikipedia number-theory item 104: Congruent number problem (a corollary to Birch and Swinnerton-Dyer conjecture, per Tunnell's theorem…}\label{statement}

\subsection{Definitions and mathematical statement}

\paragraph{Shared notation.}$\mathbb N=\{1,2,\ldots\}$, and $\mathbb Z,\mathbb Q,\mathbb R,\mathbb C$ denote the integers, rationals, reals and complex numbers. Any other conventions are specified locally.


A positive rational $q$ is congruent if there exist $a,b,c\in\mathbb Q_{>0}$ such that $a^2+b^2=c^2$ and $ab/2=q$, so that $q$ is the area of a rational right triangle. Multiplying triangle sides by $t\in\mathbb Q_{>0}$ multiplies its area by $t^2$. If $q=u/v$ is in lowest terms, then $v^2q=uv$ is an integer; removing square factors of $uv$ reduces the question to a positive squarefree integer. Equivalently, the elliptic curve $E_q:y^2=x^3-q^2x$ has a rational point with $y\ne0$.
\paragraph{Statement recorded in the supplied source.}
Classify all positive rational numbers $q$ which are congruent, giving a necessary and sufficient characterization of the existence of the displayed rational right triangle (equivalently, of a rational point on $E_q$ with nonzero $y$). \sref{statement}{2}
\subsection{Short English statement}
Determine precisely which positive rational areas occur for right triangles with all three sides rational.
\subsection{Sources}
\begin{description}
\item[\hypertarget{source:statement:1}{S1}] ULAM.AI and the UnsolvedMath Contributors, UnsolvedMath: A Curated Collection of Open Mathematics Problems, catalogue entry. CC BY 4.0. \url{https://huggingface.co/datasets/ulamai/UnsolvedMath}.
\item[\hypertarget{source:statement:2}{S2}] Chandan Singh Dalawat, \emph{Congruent numbers, elliptic curves, and the passage from the local to the global: An update}, arXiv:2201.11071v2 (2022) Section on the sign of the functional equation, p. 2, recalling Proposition 6. \url{https://arxiv.org/pdf/2201.11071}.
\item[\hypertarget{source:statement:3}{S3}] Wikipedia contributors, \emph{List of unsolved problems in mathematics}, revision 1366636767 (2026) Number theory, Wikipedia number-theory item 104: Congruent number problem (a corollary to Birch and Swinnerton-Dyer conjecture, per Tunnell's theorem…. \url{https://en.wikipedia.org/w/index.php?title=List_of_unsolved_problems_in_mathematics&oldid=1366636767#Number_theory}.
\end{description}
\label{end:statement}
\end{document}
